\section*{Capítulo \ref{ch:b2:charges-currents-conduction} --- Cargas, corrientes y conducción}

\begin{solution}{exo:b2:charges-currents-conduction:1}
$j = 16/2.5 \times 10^{-6} = \qty{6.4e6}{A/m^2}$; $v = j/ne = \qty{0.47}{mm/s}$; \qty{20}{m}
en $4.3 \times 10^4$ s, doce horas. $E = j/\gamma = \qty{0.11}{V/m}$; $U = \qty{2.1}{V}$;
$P = UI = \qty{34}{W}$.
\end{solution}

\begin{solution}{exo:b2:charges-currents-conduction:2}
Rayo, $3 \times 10^4/\pi10^{-4} = \qty{1e8}{A/m^2}$; haz, $\qty{4e3}{A/m^2}$; axón,
$10^{-9}/1.3 \times 10^{-11} = \qty{80}{A/m^2}$; corriente anular, \qty{1e-8}{A/m^2}. Solo el
rayo supera los $6 \times 10^6$ del hilo.
\end{solution}

\begin{solution}{exo:b2:charges-currents-conduction:3}
$R = \rho_eL/S$: $0.68$, $1.1$, $4.0$, \qty{44}{\ohm}. Aluminio: $2.5 \times 2.8/1.7 =
\qty{4.1}{mm^2}$, cociente de masas $(4.1 \times 2700)/(2.5 \times 8900) = 0.50$ --- la mitad de masa.
Calefactor: $R = 230^2/1000 = \qty{53}{\ohm}$, $S = \qty{0.196}{mm^2}$, $L = RS/\rho_e = \qty{9.4}{m}$.
\end{solution}

\begin{solution}{exo:b2:charges-currents-conduction:4}
$\tau = m\gamma/ne^2$: cobre \qty{2.5e-14}{s}, $\mu = e\tau/m = \qty{4.4e-3}{m^2/(V.s)}$,
$\ell = \qty{40}{nm}$; plata \qty{3.8e-14}{s}, \qty{6.7e-3}{m^2/(V.s)}, \qty{61}{nm}.
Silicio: $\gamma = ne\mu = 10^{22} \times 1.6 \times 10^{-19} \times 0.14 = \qty{220}{S/m}$; $v =
\mu E = \qty{140}{m/s}$.
\end{solution}

\begin{solution}{exo:b2:charges-currents-conduction:5}
(a) Entra corriente y no sale ninguna: $\dd Q/\dd t = I$. (b) La carga se acumula: $\partial_t
\rho \ne 0$. (c) $\vect j = (I/4\pi r^2)\vect e_r$: $\operatorname{div}\vect j = \frac1{r^2}
\partial_r(I/4\pi) = 0$. (d) Solo en el electrodo, donde el hilo inyecta la
corriente --- una fuente del flujo.
\end{solution}

\begin{solution}{exo:b2:charges-currents-conduction:6}
(a) $j = I/2\pi rL$, $E = j/\gamma$, radiales. (b) $U = \int_a^bE\,\dd r = I\ln(b/a)/
2\pi\gamma L$: $R = \ln3.5/(2\pi \times 10^{-13} \times 10^3) = \qty{2e9}{\ohm}$ por kilómetro.
(c) \qty{0.5}{\micro A}, \qty{0.5}{mW}. (d) $RC = \varepsilon_0\varepsilon_r/\gamma = \qty{200}{s}$ para
$\varepsilon_r = 2.3$: el cable, cargado y aislado, se descarga a través de su
propio aislante con el tiempo de relajación de ese aislante.
\end{solution}

\begin{solution}{exo:b2:charges-currents-conduction:7}
(a) $\omega\tau$: $8 \times 10^{-12}$, $1.6 \times 10^{-3}$, $4.7$, $79$. (b) $|\underline\gamma|/\gamma_0 =
1/\sqrt{1 + \omega^2\tau^2}$: $1$, $1$, $0.21$, $0.013$; fase $-\arctan\omega\tau$: $0$,
$-\qty{0.09}{\degree}$, $-\qty{78}{\degree}$, $-\qty{89}{\degree}$. (c) $\omega\tau \le 0.14$: $f \le
\qty{0.9}{THz}$. (d) $\underline\gamma \approx ne^2/\iu m\omega$: $\underline{\vect j}$ va
\qty{90}{\degree} por detrás de $\underline{\vect E}$, $\langle\vect j\cdot\vect E\rangle = 0$: no hay
efecto Joule --- los electrones oscilan libremente, como en un plasma.
\end{solution}

\begin{solution}{exo:b2:charges-currents-conduction:8}
(a) $S = \qty{3.1e-8}{m^2}$, $j = \qty{3.2e8}{A/m^2}$, $p = j^2/\gamma = \qty{1.7e9}{W/m^3}$.
(b) Energía necesaria para fundirlo por unidad de volumen $\rho(c\Delta T + L_f) = 8900(4.1 \times 10^5 + 2.05 \times
10^5) = \qty{5.5e9}{J/m^3}$: \qty{3.2}{s}. (c) Nueve veces más rápido: \qty{0.36}{s}. (d) La
corriente nominal es aquella cuyo calentamiento por efecto Joule evacua el
aire justo por debajo del punto de fusión; la arena apaga el arco que, de
otro modo, seguiría conduciendo por el hueco tras fundirse el hilo.
\end{solution}

\begin{solution}{exo:b2:charges-currents-conduction:9}
(a) $\partial_t\rho = -\operatorname{div}\vect j = -\gamma\operatorname{div}\vect E = -\gamma\rho/
\varepsilon_0$. (b) $\varepsilon_0/\gamma$: \qty{1.5e-19}{s}, \qty{1.8e-12}{s}, \qty{1.8e-10}{s},
\qty{9}{s}, $9 \times 10^4$ s. (c) $f \lesssim 1/2\pi\tau_r$: $10^{18}$ Hz (el \omterm{prop:b2:charges-currents-conduction:drude}{modelo de Drude}
falla mucho antes), \qty{90}{GHz}, \qty{0.9}{GHz}, \qty{0.02}{Hz}, \qty{2e-6}{Hz}. (d)
Una varilla de vidrio perdería su carga en un minuto a través de su
volumen; el vidrio seco real está más cerca de $10^{-14}$ S/m (horas), y lo
que suele descargarla es la humedad superficial --- coherente en orden de
magnitud.
\end{solution}

\begin{solution}{exo:b2:charges-currents-conduction:10}
(a) $\vect j = (I/2\pi r^2)\vect e_r$, $E = I/2\pi\gamma r^2$, $V = I/2\pi\gamma r$. (b) $R = V(a)
/I = 1/2\pi\gamma a = \qty{32}{\ohm}$. (c) $V_{\text{pica}} = \qty{640}{kV}$; $V(10) = \qty{32}{kV}$,
$V(10.7) = \qty{30}{kV}$: \qty{2.1}{kV} entre los pies. (d) Las patas
delanteras y traseras distan \qty{1.5}{m} en la dirección radial: de $4$ a
\qty{5}{kV}, y el camino de la corriente pasa por el corazón.
\end{solution}

\begin{solution}{exo:b2:charges-currents-conduction:11}
(a) $E_H = jB/nq$ con $j = I/wt$: $V_H = E_Hw = IB/nqt$. (b) $1/(8.5 \times 10^{28}
\times 1.6 \times 10^{-19} \times 10^{-4}) = \qty{0.7}{\micro V}$. (c) \qty{6.3}{mV}: $V_H \propto 1/n$,
y los semiconductores tienen $10^6$ veces menos portadores. (d) $B = V_Hnqt/I =
10^{-5} \times 10^{22} \times 1.6 \times 10^{-19} \times 2 \times 10^{-5}/10^{-3} = \qty{0.32}{T}$; el signo
de $V_H$ da el signo de los portadores (para un $\vect B$ conocido) o la
dirección de $\vect B$ (para un sensor conocido).
\end{solution}

\begin{solution}{exo:b2:charges-currents-conduction:12}
(a) $\vect j = ne\mu_+\vect E + (-ne)(-\mu_-\vect E)$: ambas especies se
desplazan en sentidos opuestos y llevan corriente en el mismo sentido. (b) $n =
0.5 \times 10^3 \times 6.02 \times 10^{23} = \qty{3e26}{m^{-3}}$: $\gamma = 3 \times 10^{26} \times 1.6 \times 10^{-19}
\times 1.31 \times 10^{-7} = \qty{6.3}{S/m}$ (a esta concentración los iones se
estorban entre sí: \qty{5}{S/m} medidos). (c) \qty{0.52}{\micro m/s} y \qty{0.79}{\micro m/s}:
el Cl$^-$ lleva el $60\%$. (d) $R = L/\gamma S = \qty{200}{\ohm}$, $U = \qty{200}{V}$, $P =
\qty{200}{W}$ sobre \qty{10}{g} de agua: \qty{10}{K} en \qty{2}{s} --- de ahí las
placas de refrigeración de las cubetas de electroforesis.
\end{solution}

\begin{solution}{pb:b2:charges-currents-conduction:1}
\textbf{1.} $n = \rho N_A/M = \qty{8.4e28}{m^{-3}}$.

\textbf{2.} $j = \qty{6.4e6}{A/m^2}$, $v = j/ne = \qty{0.47}{mm/s}$; \qty{35}{min} por metro.

\textbf{3.} $\tau = m\gamma/ne^2 = \qty{2.5e-14}{s}$, $\mu = \qty{4.4e-3}{m^2/(V.s)}$, $\ell =
\qty{40}{nm}$, $150$ distancias atómicas: los electrones chocan con defectos
y vibraciones térmicas, no con los átomos.

\textbf{4.} $\rho_e \propto T$ da $\alpha = 1/293 = \qty{3.4e-3}{K^{-1}}$, cercano al
$4 \times 10^{-3}$ medido.

\textbf{5.} $E = \qty{0.11}{V/m}$, $U = \qty{0.11}{V}$; $\tfrac12mv^2 = \qty{1e-37}{J}$ frente a
$k_BT = \qty{4e-21}{J}$: el gas de electrones no nota la corriente.

\textbf{6.} $I/e = \qty{1e20}{}$ electrones por segundo.

\textbf{7.} $Q = neSL = \qty{3.4e4}{C}$; $QE = \qty{3.6}{kN}$, frente a un peso de
\qty{0.2}{N}. El campo tira de la red positiva con la fuerza opuesta; los
electrones ceden su cantidad de movimiento a la red por colisiones: el
hilo en conjunto no siente nada.

\textbf{8.} $p = j^2/\gamma = \qty{6.8e5}{W/m^3}$, \qty{1.7}{W/m}.

\textbf{9.} Superficie $\pi d = \qty{5.6e-3}{m^2}$ por metro: $\Delta T = 1.7/(10 \times 5.6 \times
10^{-3}) = \qty{30}{K}$; la resistividad sube un $12\%$ y la potencia a \qty{1.9}{W}:
$\Delta T \approx \qty{35}{K}$.

\textbf{10.} $\Delta T = 1.7(1 + 0.004\Delta T)/(2 \times 5.6 \times 10^{-3})$: $\Delta T(1 - 0.6) = 150$,
$\Delta T \approx \qty{400}{K}$ --- cerca del embalamiento térmico: los cables
enterrados deben usarse por debajo de su corriente nominal.

\textbf{11.} $p/\rho c = \qty{0.2}{K/s}$; a \qty{160}{A}, \qty{20}{K/s}: se funde
en un minuto aproximadamente (menos, según sube $\rho_e$); el
interruptor automático abre en milisegundos.

\textbf{12.} $\omega\tau = 8 \times 10^{-12}$: la ley de Ohm no se altera.

\textbf{13.} $\varepsilon_0/\gamma = \qty{1.5e-19}{s}$: ninguna carga en el interior; el
campo del hilo lo crean cargas de su \emph{superficie}, cuya densidad
varía a lo largo del hilo.

\textbf{14.} $Q = I_0\tau_c = \qty{16}{mC}$, $V = Q/C = \qty{160}{V}$.

\textbf{15.} El lado del hilo lleva $I$ y el lado del hueco no lleva
corriente de conducción: la carga crece, y con ella el campo eléctrico del
hueco.

\textbf{16.} $I = 16\eu^{-0.5} = \qty{9.7}{A}$; $E = Q/\varepsilon_0A$, $\dd E/\dd t = I/\varepsilon_0A
= 9.7/(8.85 \times 10^{-12} \times 10^{-4}) = \qty{1.1e16}{V.m^{-1}.s^{-1}}$.

\textbf{17.} $\varepsilon_0A\,\dd E/\dd t = I$: la corriente de desplazamiento.

\textbf{18.} $j = \gamma E = 10^{-14} \times 2.3 \times 10^5 = \qty{2e-9}{A/m^2}$, $10^{-14}$ A
por la sección: nada; ruptura a \qty{3}{kV} en el milímetro.

\textbf{19.} $R = 1000/(6 \times 10^7 \times 10^{-5}) = \qty{1.7}{\ohm}$; $\Delta U = \qty{67}{V}$; $P =
\qty{2.7}{kW}$, el $29\%$ de los \qty{9.2}{kW} entregados.

\textbf{20.} $I = \qty{0.46}{A}$, $\Delta U = \qty{0.8}{V}$, \qty{0.35}{W}: transportar a
alta tensión divide la pérdida por $(U_2/U_1)^2$.

\textbf{21.} $10 \times 6/3.6 = \qty{17}{mm^2}$; \qty{89}{kg/km} de cobre frente a
\qty{45}{kg/km} de aluminio: más ligero y más barato en las torres; en la
pared, la sección menor del cobre cabe en los tubos y en los bornes.

\textbf{22.} $V_H = IB/nqt = 5 \times 10^{-3} \times 0.1/(10^{22} \times 1.6 \times 10^{-19} \times 10^{-4}) =
\qty{3.1}{mV}$.

\textbf{23.} $+24\%$: \qty{2.1}{\ohm}, \qty{3.3}{kW} perdidos.

\textbf{24.} $RI^2t = 1.7 \times 4 \times 10^8 \times 10^{-4} = \qty{67}{kJ}$ sobre \qty{89}{kg}: \qty{2}{K}.
$j = \qty{2e9}{A/m^2}$, trescientas veces el valor doméstico.

\textbf{25.} $n$ (la materia) y $\tau$ (las colisiones) dan $\gamma = ne^2\tau/m$;
$\gamma$ enlaza $j$ y $E$ (Ohm); $j$ se integra hasta $I$ y $E$ hasta $U$ (el
circuito); $j^2/\gamma$ es el calor.
\end{solution}
